% Modernised reading — fonds Alexandre Grothendieck, shelfmark 106, pages 1-16.
%
% This is an INTERPRETATION, not a transcription. Everything the critical
% apparatus carried moves into footnotes. In any dispute of substance the
% transcription governs: whoever cites Grothendieck must cite batch-01.fr.tex.
%
% Pass: Opus 5 (claude-opus-5), 2026-09-13 — first pass, unchecked against the
% pages by a human. The folder was read whole; it holds one batch, pages 1-16.
% His notebook leaves are pages 1, 3, 5, 7, 9, 11, 13, 15 (his 1 to 8); pages
% 2, 12 and 14 are computer listings carrying trial figures of his; pages 4,
% 6, 8, 10 and 16 carry nothing of his.
%
% What the folder is. One argument about the homotopy category of a category
% with weak equivalences W and cofibrations cof satisfying Baues's axioms
% C1-C3. First (pages 1-5): the inclusion of cofibrant objects induces an
% equivalence of localisations, reduced to a condition (*). Then (pages 5-15),
% every object being cofibrant: an attempt to compute W^{-1}C by a calculus of
% fractions x -> y~ <- y, which breaks twice — the inverse of a weak
% equivalence does not come out (page 9), and faithfulness is not clear (page
% 11) — is rectified by changing the definition, and ends on a criterion (**)
% for faithfulness, stated and not proved.
%
% Notation fixed for the whole document. C, W, cof; C_cof the full
% subcategory of cofibrant objects, W_cof = W restricted to it; Ho = W^{-1}C.
% « cof » kept, as he writes it. A tilde under an arrow marks membership of W
% and is set as a label. x~ -> x (alpha_x) the cofibrant replacement. Two
% homonyms separated here: Ch_0(x,y) is page 5's category (i in cof ∩ W, f
% arbitrary) and Ch'_0(x,y) is page 11's rectified one (f in cof, i in W);
% C~ and C~' the corresponding categories of components. He writes Ch_0 and
% C~ for both.
%
% Substantive departures from the pages, each footnoted where it occurs:
% — page 5: the composition law (g,j) o (f,i) = (g'g) . k_y omits f; the
%   composite is (g'f, kj), g' and k coming from the pushout of g along i.
% — page 7: f |-> (f, 1d_x) must be (f, 1_y).
% — page 11: his doubt that inverting W_0 = W ∩ cof suffices is unfounded
%   when every object is cofibrant: the factorisation of page 9 gives
%   f = alpha i with alpha g = 1_y and g in W ∩ cof, so [alpha] = [g]^{-1}
%   and [f] are invertible in W_0^{-1}C (K. S. Brown's factorisation lemma).
%   The reading says so; what remains open on the leaves is faithfulness.
% — page 1: the « weak lifting lemma » is used in a form the reading derives
%   from C1-C3 (factor (alpha_y, f alpha_x) out of y~ v x~).
%
% Modern names given and footnoted as ours: cofibration category (Baues),
% Brown's factorisation lemma, calculus of fractions (Gabriel-Zisman),
% homotopy calculus of fractions, cofibrant replacement.
%
% What the folder announces and never establishes: the independence of the
% choices in the rectified construction, reduced to (**); the sufficiency of
% (**) for faithfulness; hence that C~' -> W^{-1}C is an isomorphism. The
% illegible marginal note of page 13 is not restored.
\documentclass[11pt,a4paper]{article}
% Le préambule commun à toutes les transcriptions du fonds Grothendieck.
%
% Il tient en une idée : **l'incertitude se compose, elle ne se résout pas.**
% Une transcription qui lisse un mot douteux a détruit la seule chose pour
% laquelle on la faisait — un lecteur doit pouvoir savoir, à chaque endroit,
% ce qui était sur le feuillet et ce qui a été ajouté. Les six macros qui
% suivent sont donc l'appareil critique, et non de la décoration.
%
% Les mêmes commandes sont reconnues par `scripts/render.mjs`, qui produit la
% vue de lecture du site. Une macro ajoutée ici doit l'être aussi là-bas,
% faute de quoi le rendu échouera bruyamment — ce qui est voulu : un rendu
% silencieusement faux est pire qu'un rendu absent.

\ProvidesPackage{grothendieck}[2026/08/08 Transcriptions du fonds Grothendieck]

\usepackage{amsmath,amssymb,amsthm}
\usepackage{mathrsfs}
\usepackage{stmaryrd}
% Les diagrammes commutatifs sont partout dans ce fonds : c'est la langue
% même de la géométrie algébrique de Grothendieck, et une transcription qui
% les rendrait en prose ne transcrirait rien.
\usepackage{tikz-cd}
% Français par défaut — c'est la langue des feuillets — mais l'anglais est
% chargé aussi : la traduction et le résumé sont des documents anglais, et la
% typographie française (espaces avant les deux-points) y serait une faute.
% Ces documents basculent par \selectlanguage{english} en tête de corps.
\usepackage[english,french]{babel}
\usepackage{fontspec}
\usepackage{geometry}
\geometry{a4paper,margin=25mm,marginparwidth=28mm}
\usepackage{marginnote}
\usepackage{xcolor}
% ulem, not soul. `soul`'s \st and \ul are built on a character-by-character
% scan that XeLaTeX's font machinery breaks: the document fails with an
% undefined control sequence rather than degrading. ulem does the same two jobs
% and is Unicode-engine safe. `normalem` stops it hijacking \emph, which the
% transcriptions use for Grothendieck's own emphasis.
\usepackage[normalem]{ulem}
\usepackage{hyperref}
\hypersetup{colorlinks=true,linkcolor=black,urlcolor=black,pdfborder={0 0 0}}

% --- Métadonnées du lot -----------------------------------------------------
% Reprises telles quelles par le script de rendu, qui les affiche en tête de
% la vue de lecture. Elles fixent ce que le fichier prétend couvrir : sans
% elles, une transcription flotte sans point d'attache dans un fonds de seize
% mille feuillets.

\newcommand{\folder}[1]{\gdef\grothFolder{#1}}
\newcommand{\batch}[1]{\gdef\grothBatch{#1}}
\newcommand{\leaves}[2]{\gdef\grothFirst{#1}\gdef\grothLast{#2}}
\newcommand{\foldertitle}[1]{\gdef\grothFoldertitle{#1}}
\gdef\grothFolder{?}\gdef\grothBatch{?}\gdef\grothFirst{?}\gdef\grothLast{?}\gdef\grothFoldertitle{}

% --- Le résumé --------------------------------------------------------------
%
% Il ouvre la lecture modernisée, sous le titre « Résumé », et il est composé
% en petit. Ce n'est pas un ornement : c'est le seul passage du document qui
% ne soit pas des mathématiques, et un lecteur doit pouvoir le reconnaître
% avant de l'avoir lu — puis le sauter s'il connaît déjà le sujet, ou s'y
% arrêter s'il ne le connaît pas. Le filet qui le suit marque la reprise du
% corps.

\newenvironment{resume}
  {\small\setlength{\parskip}{0.45em}}
  {\par\medskip\hrule\medskip}

% --- Mots-clés ---------------------------------------------------------------
%
% En anglais, en clôture du résumé de la lecture modernisée : le vocabulaire
% moderne sous lequel chercher ce que les feuillets construisent. Le manifeste
% du site (`npm run manifest`) les extrait pour étiqueter la cote entière —
% cette ligne est la source unique des tags, il n'y a pas de fichier à côté.
\newcommand{\keywords}[1]{\par\smallskip\noindent
  {\footnotesize\textsc{Keywords} --- \textit{#1}}\par}

% --- L'appareil critique ----------------------------------------------------

% Le numéro de feuillet, en marge. C'est l'ancre qui permet de régler un
% désaccord sur une formule en regardant *un* feuillet plutôt qu'une chemise
% de six cents. Le site s'en sert aussi pour tourner les pages du fac-similé
% quand on fait défiler la transcription.
\newcommand{\leaf}[1]{%
  \marginnote{\footnotesize\color{gray!70}\textsf{#1}}%
  \label{leaf:#1}\ignorespaces}

% Illisible : on ne devine pas. Un mot inventé qui se lit comme les autres est
% le pire résultat possible.
\newcommand{\ill}{\textcolor{red!60!black}{[\,\ldots\,]}}

% Lecture douteuse : le mot est donné, et signalé comme incertain.
\newcommand{\uncertain}[1]{\uline{#1}}

% Ajout de l'éditeur — une lettre manquante, un mot sous-entendu.
\newcommand{\add}[1]{\textcolor{blue!50!black}{[#1]}}

% Biffé par Grothendieck lui-même. Ce qu'il a rayé fait partie du document :
% une rature dit souvent où la pensée a changé de direction.
\newcommand{\struck}[1]{\sout{#1}}

% Note du transcripteur, sur l'état matériel du feuillet.
\newcommand{\note}[1]{\par\smallskip{\small\color{gray!60!black}\textit{[#1]}}\par\smallskip}

% Note marginale de Grothendieck, distincte du corps du texte.
\newcommand{\marginal}[1]{\par\smallskip\noindent\hspace{1em}%
  {\small\color{gray!60!black}#1}\par\smallskip}

% --- Titre ------------------------------------------------------------------

\AtBeginDocument{%
  \begin{center}
    {\large\bfseries Fonds Alexandre Grothendieck --- cote n\textsuperscript{o}~\grothFolder}\\[2pt]
    {\itshape\grothFoldertitle}\\[4pt]
    {\small feuillets \grothFirst--\grothLast\ (lot \grothBatch)}\\[2pt]
    {\footnotesize\color{gray!60!black}Transcription établie d'après le fac-similé numérisé
     par l'Université de Montpellier.\\
     Les lectures incertaines sont soulignées, les illisibles notées [\,\ldots\,],
     les ajouts entre crochets.}
  \end{center}
  \vspace{1em}}

\endinput


\folder{106}
\pages{1}{16}
\dating{[à partir de 1982]}
\watermark{Édition de démonstration\\interprétation personnelle de l'œuvre}
\foldertitle{[Champs (stacks) 3] : notes manuscrites (s.d.) — lecture modernisée du dossier entier}

\begin{document}

\section*{Résumé}

\begin{resume}
En topologie, deux espaces qu'on peut déformer continûment l'un dans l'autre
sont considérés comme « le même » du point de vue de l'homotopie. Pour en
faire une mathématique, on part d'une catégorie d'objets et d'une classe de
flèches qu'on déclare être des équivalences --- les \emph{équivalences
faibles} --- et l'on fabrique la catégorie où ces flèches sont devenues
inversibles : la \emph{catégorie homotopique}. La construction formelle est
simple à écrire et presque impossible à utiliser, car une flèche y est un
zigzag arbitrairement long de flèches et d'inverses formels.

Tout l'art est de raccourcir les zigzags. Si l'on dispose en plus d'une
classe de \emph{cofibrations} --- les « bonnes » inclusions, le long
desquelles on peut recoller --- avec quelques axiomes, on espère que toute
flèche de la catégorie homotopique s'écrive comme un aller simple suivi d'un
retour inversible, $x \to \widetilde{y} \leftarrow y$, exactement comme une
fraction $a/s$ s'écrit avec un seul dénominateur. C'est ce qu'on appelle un
calcul de fractions.

Ces huit feuillets de bloc, écrits en août 1982 ou après, font cet essai sous
les axiomes de H.-J.~Baues. Ils commencent par un point acquis : on peut se
limiter aux objets « cofibrants », sans rien perdre. Puis ils tentent le
calcul de fractions, et l'on voit l'essai échouer à deux reprises, en clair
--- « Ça se présente mal ! », « je doute qu'il suffise… » ---, être rectifié
par un changement de définition, et se réduire pour finir à une condition
précise sur les sections d'une équivalence faible, qui n'est pas démontrée.

C'est une page de travail au sens propre : on y voit l'endroit exact où une
construction résiste, et la manière dont on déplace la définition plutôt que
de forcer la preuve. L'un des deux doutes est d'ailleurs levé par les feuillets
eux-mêmes, deux pages plus tôt, sans qu'il le remarque.

\keywords{cofibration category, localization, homotopy category, cofibrant replacement, Brown factorization lemma, calculus of fractions}
\end{resume}

\section*{Le fil du dossier, et les conventions}

\begin{enumerate}
\item[1.] \emph{Objets cofibrants} (pages 1 à 5) :
$W_{\mathrm{cof}}^{-1}\mathcal{C}_{\mathrm{cof}} \to W^{-1}\mathcal{C}$ est une
équivalence si et seulement si une condition $(\star)$ est vérifiée.
\item[2.] \emph{Premier calcul de fractions} (pages 5 à 9) : les catégories
$\mathrm{Ch}_{0}(x,y)$, leur composition, la catégorie
$\widetilde{\mathcal{C}}$, et l'échec sur les inverses.
\item[3.] \emph{Le foncteur $q$} (page 9) : bijectif sur les objets, plein.
\item[4.] \emph{Rectification} (pages 11 à 15) : la définition
$\mathrm{Ch}'_{0}(x,y)$, la comparaison de factorisations, et le critère
$(\star\star)$.
\end{enumerate}

\emph{Conventions.} $\mathcal{C}$ est une catégorie munie de deux classes de
flèches, les équivalences faibles $W$ et les cofibrations
$\mathrm{cof}$, satisfaisant aux axiomes C1 à C3 de Baues, avec un objet
initial $\emptyset$.\footnote{Ce sont les axiomes d'une \emph{catégorie à
cofibrations} au sens de H.-J.~Baues (exposés dans \emph{Algebraic Homotopy},
1989) : C1, les isomorphismes sont dans $W$ et dans $\mathrm{cof}$, et $W$
satisfait au deux-sur-trois ; C2, les sommes amalgamées le long d'une
cofibration existent, et les cofibrations, ainsi que les cofibrations
triviales, sont stables par ce changement de base ; C3, toute flèche se factorise en une
cofibration suivie d'une équivalence faible. Il écrit « C1 : C3 de Baues » et
ne les rappelle pas ; la page 5 dit que les seules sommes amalgamées « dont
on postule l'existence » sont celles le long d'une cofibration. Le listing
des pages 12 et 14 est daté par la machine des 24 et 25 août 1982.}
$\mathcal{C}_{\mathrm{cof}}$ est la sous-catégorie pleine des objets
cofibrants et $W_{\mathrm{cof}}$ la restriction de $W$. Une flèche marquée
$\sim$ est dans $W$.

\emph{Deux définitions sous un même nom.} La page 5 appelle
$\underline{Ch}_{0}(x,y)$ la catégorie des diagrammes
$x \xrightarrow{f} \widetilde{y} \xleftarrow{i} y$ avec $i \in \mathrm{cof}
\cap W$ ; la page 11 reprend le même nom pour les diagrammes où
$f \in \mathrm{cof}$ et $i \in W$. On écrit ici $\mathrm{Ch}_{0}$ pour la
première, $\mathrm{Ch}'_{0}$ pour la seconde, et $\widetilde{\mathcal{C}}$,
$\widetilde{\mathcal{C}}'$ pour les catégories qu'on en tire.

\emph{Ce que le dossier annonce et n'établit pas} : que le choix des
factorisations dans la construction rectifiée n'importe pas, ramené à
$(\star\star)$ ; que $(\star\star)$ suffise à la fidélité ; et donc que
$\widetilde{\mathcal{C}}' \to W^{-1}\mathcal{C}$ soit un isomorphisme.

\pagerange{1}{3}
\section*{Les objets cofibrants suffisent (pages 1 à 3)}

\textbf{Énoncé connu.} Le foncteur $\Psi :
W_{\mathrm{cof}}^{-1}\mathcal{C}_{\mathrm{cof}} \to W^{-1}\mathcal{C}$ induit
par l'inclusion est une équivalence de catégories.\footnote{Il écrit « Alors il
est connu que », entre parenthèses. L'énoncé est en effet classique pour les
catégories à cofibrations. Une note au coin de la page, coupée d'un trait,
porte « Ho $\mathcal{C}$ » et « cat.\ de […] co/fibrations ? », l'abréviation
du milieu ne se résolvant pas.} Pour le voir, on construit un quasi-inverse
$\Phi : W^{-1}\mathcal{C} \to W_{\mathrm{cof}}^{-1}\mathcal{C}_{\mathrm{cof}}$,
c'est-à-dire un foncteur $\mathcal{C} \to
W_{\mathrm{cof}}^{-1}\mathcal{C}_{\mathrm{cof}}$ qui rend inversibles les
flèches de $W$.

\emph{Remplacement cofibrant.} Pour tout $x$, C3 appliqué à $\emptyset \to x$
donne $\alpha_{x} : \widetilde{x} \to x$ dans $W$ avec $\widetilde{x}$
cofibrant.

\emph{Relèvement faible.} Soit $f : x \to y$. On ne peut pas relever
$f\alpha_{x} : \widetilde{x} \to y$ à $\widetilde{y}$, mais on peut le relever à
un $\widetilde{y}'$ : il existe $i : \widetilde{y} \to \widetilde{y}'$ dans
$\mathrm{cof} \cap W$, $\alpha'_{y} : \widetilde{y}' \to y$ dans $W$ avec
$\alpha_{y} = \alpha'_{y}\, i$, et $\widetilde{f} : \widetilde{x} \to
\widetilde{y}'$ avec $\alpha'_{y}\widetilde{f} = f\alpha_{x}$.

\begin{tikzcd}[column sep=small, row sep=small, nodes={font=\scriptsize}]
\widetilde{x} \arrow[r, "\widetilde{f}"] \arrow[d, "\underset{\sim}{\alpha_x}"'] & \widetilde{y}\,' \arrow[d, "\underset{\sim}{\alpha'_y}"] & \widetilde{y} \arrow[l, "i"'] \arrow[dl, "\underset{\sim}{\alpha_y}"] \\
x \arrow[r, "f"'] & y &
\end{tikzcd}

On l'obtient en factorisant $(\alpha_{y}, f\alpha_{x}) : \widetilde{y} \vee
\widetilde{x} \to y$ en une cofibration suivie d'une équivalence faible ;
$\widetilde{y} \to \widetilde{y} \vee \widetilde{x}$ est une cofibration parce
que $\widetilde{x}$ est cofibrant, et $i$ est dans $W$ par
deux-sur-trois.\footnote{La page invoque le « weak lifting lemma » de Baues
sans l'énoncer ; on donne la forme sous laquelle elle l'utilise et la manière
de l'obtenir sous C1--C3.}

On pose $\Phi(f) = [i]^{-1}[\widetilde{f}] : \widetilde{x} \to \widetilde{y}$
dans $W_{\mathrm{cof}}^{-1}\mathcal{C}_{\mathrm{cof}}$.

\emph{Indépendance des choix.} Deux choix $(i_{1}, \alpha'_{1},
\widetilde{f}_{1})$ et $(i_{2}, \alpha'_{2}, \widetilde{f}_{2})$ se comparent
dans la somme amalgamée $\widetilde{y}'_{1} \vee_{\widetilde{y}}
\widetilde{y}'_{2}$, et l'on est ramené à la condition :

\begin{enumerate}
\item[$(\star)$] si $x', y'$ sont cofibrants, $\varphi, \psi : x' \to y'$, et
$\alpha : y' \to y$ est dans $W$ avec $\alpha\varphi = \alpha\psi$, alors
$[\varphi] = [\psi]$ dans $W_{\mathrm{cof}}^{-1}\mathcal{C}_{\mathrm{cof}}$.
\end{enumerate}

Sous $(\star)$, on obtient un foncteur $\mathcal{C} \to
W_{\mathrm{cof}}^{-1}\mathcal{C}_{\mathrm{cof}}$, la compatibilité aux
identités et aux compositions étant immédiate.\footnote{Le feuillet de listing
de la page 2, tourné d'un quart de tour, est couvert d'essais de ce même
diagramme --- deux relèvements dans $\widetilde{y}'$ et $\widetilde{y}''$ et leur
somme amalgamée sous $\widetilde{y}$ ---, sans étiquettes.}

\emph{Il rend $W$ inversible} (page 3). Si $f \in W$, alors $f\alpha_{x}$ et
$\alpha'_{y}$ sont dans $W$, donc $\widetilde{f}$ aussi, et $[\widetilde{f}]$
est un isomorphisme.

\emph{Les deux composés.} $\Phi\Psi$ est l'identité : pour $x$ cofibrant on
choisit $\widetilde{x} = x$ et $\alpha_{x} = 1$. Pour $\Psi\Phi$, il faut voir
que $R : x \mapsto \widetilde{x}$, $\mathcal{C} \to W^{-1}\mathcal{C}$, est
isomorphe au foncteur canonique par les isomorphismes $[\alpha_{x}]$,
c'est-à-dire que
\[
[\alpha_{y}]\,\bigl([i_{f}]^{-1}[\widetilde{f}]\bigr) = [f]\,[\alpha_{x}] .
\]
Or $\alpha_{y} = \beta_{f}\, i_{f}$, où $\beta_{f} = \alpha'_{y}$, donne
$[\alpha_{y}][i_{f}]^{-1} = [\beta_{f}]$, et $[\beta_{f}\widetilde{f}] =
[f\alpha_{x}] = [f][\alpha_{x}]$.

\pagerange{5}{5}
\section*{La proposition, et le passage aux objets cofibrants (page 5)}

\textbf{Proposition.} Pour que $\Psi :
W_{\mathrm{cof}}^{-1}\mathcal{C}_{\mathrm{cof}} \to W^{-1}\mathcal{C}$ soit une
équivalence de catégories, il faut et il suffit que la condition $(\star)$ soit
vérifiée ; elle est d'ailleurs impliquée par la fidélité de $\Psi$.

La suffisance est ce qui précède. La nécessité vient de ce que $[\varphi] =
[\psi]$ dans $W^{-1}\mathcal{C}$, $[\alpha]$ y étant inversible.\footnote{Il
ajoute un « ? » en interligne après « il f.\ et s. ». Comme l'énoncé
initial est vrai pour les catégories à cofibrations, $(\star)$ l'est aussi : la
proposition ne dit rien de faux, et le point d'interrogation porte sur ce que
les feuillets démontrent, non sur ce qui est vrai.}

\emph{Réduction.} On travaille désormais dans $\mathcal{C}_{\mathrm{cof}}$, en
changeant les notations : \emph{tout objet de $\mathcal{C}$ est cofibrant}.
C'est légitime, car $\mathcal{C}_{\mathrm{cof}}$, avec les classes induites,
satisfait encore C1--C3, l'inclusion commutant aux sommes amalgamées le long
d'une cofibration --- « les seules sommes dont on postule l'existence ! ».

\pagerange{5}{7}
\section*{Premier calcul de fractions (pages 5 et 7)}

Pour deux objets $x, y$, soit $\mathrm{Ch}_{0}(x,y)$ la petite catégorie des
diagrammes
\[
x \xrightarrow{\;f\;} \widetilde{y} \xleftarrow{\;i\;} y,
\qquad i \in \mathrm{cof} \cap W,
\]
une flèche $(f, i) \to (f', i')$ étant une flèche $u : \widetilde{y} \to
\widetilde{y}'$ telle que $uf = f'$ et $ui = i'$.\footnote{La page donne
seulement les objets, et dit la catégorie « petite, connexe » suivi d'un mot
illisible ; les flèches sont celles que le passage aux composantes connexes
de la page 7 suppose.}

\emph{Composition.} Pour $(f, i) \in \mathrm{Ch}_{0}(x,y)$ et $(g, j) \in
\mathrm{Ch}_{0}(y,z)$, soit $\widetilde{z}'$ la somme amalgamée de
$\widetilde{y} \xleftarrow{i} y \xrightarrow{g} \widetilde{z}$, avec
$g' : \widetilde{y} \to \widetilde{z}'$ et $k : \widetilde{z} \to
\widetilde{z}'$. Comme $i \in \mathrm{cof} \cap W$, $k$ l'est aussi, par C2, et
l'on pose
\[
(g, j) \circ (f, i) = (g' f,\ k j) \in \mathrm{Ch}_{0}(x,z).
\]

\begin{tikzcd}[column sep=small, row sep=small, nodes={font=\scriptsize}]
x \arrow[r, "f"] & \widetilde{y} \arrow[d, "g'"'] & y \arrow[l, "i"'] \arrow[d, "g"] & \\
& \widetilde{z}\,' & \widetilde{z} \arrow[l, "k"] & z \arrow[l, "j"']
\end{tikzcd}

Le composé est bien défini à isomorphisme près, la somme amalgamée
l'étant.\footnote{La page écrit « $(g,j) \circ (f,i) = (g'g) \cdot k_{y}$ »,
dont le membre de droite ne fait pas intervenir $f$ ; le dernier facteur est
lu $k_{y}$ avec doute. Le composé que la somme amalgamée fournit est celui
qu'on écrit. Sur le feuillet, $k$ est marqué « cocart » et $g'$ désigne
bien le composé $\widetilde{y} \to \widetilde{z}'$.}

\emph{La catégorie $\widetilde{\mathcal{C}}$} (page 7). Cette composition est
associative à isomorphisme canonique près et admet pour unités, à
isomorphisme canonique près, les $\underline{1}_{x} = (1_{x}, 1_{x})$. On a
des foncteurs $\mathrm{Hom}_{\mathcal{C}}(x,y)_{\text{discret}} \to
\mathrm{Ch}_{0}(x,y)$, $f \mapsto (f, 1_{y})$, compatibles aux
compositions.\footnote{La page écrit $(f, 1d_{x})$, avec un chiffre un ; la
seconde composante doit être une flèche $y \to y$.} Posant
\[
\pi_{0}(x,y) = \pi_{0}\bigl(\mathrm{Ch}_{0}(x,y)\bigr),
\]
la composition passe aux composantes connexes et y devient associative et
unitaire. D'où une catégorie $\widetilde{\mathcal{C}}$, de mêmes objets que
$\mathcal{C}$, et un foncteur $\mathcal{C} \to \widetilde{\mathcal{C}}$.

\emph{La question.} Ce foncteur rend-il inversibles les flèches de $W$ ? Pour
$f : x \to y$ dans $W$, on cherche $(g, i)$ tel que $(g,i) \circ (f, 1_{y})$
soit dans la composante de $\underline{1}_{x}$ et $(f, 1_{y}) \circ (g,i)$ dans
celle de $\underline{1}_{y}$.

\pagerange{9}{9}
\section*{L'échec, et le foncteur $q$ (page 9)}

\emph{La factorisation.} Soit $f : x \to y$ dans $W$. Factorisons
$(f, 1_{y}) : x \sqcup y \to y$ en une cofibration $x \sqcup y \to \bar{x}$
suivie d'une équivalence faible $\alpha : \bar{x} \to y$, et notons $i$ et $g$
les composés de $x \to x \sqcup y$ et $y \to x \sqcup y$ avec
$x \sqcup y \to \bar{x}$. Alors
\[
\alpha i = f, \qquad \alpha g = 1_{y}, \qquad i, g \in \mathrm{cof} \cap W,
\]
$i$ et $g$ étant des cofibrations parce que $x$ et $y$ sont cofibrants, et des
équivalences faibles par deux-sur-trois.\footnote{La page présente la
factorisation par le « weak lifting lemme », avec un diagramme pointillé, et
note en marge « on utilise $y \in \mathcal{C}_{\mathrm{cof}}$ ici ». Elle dit
$i \in W \cap \mathrm{cof}$ et « $\alpha$ admet une section $g$ », sans
relever que $g$ est aussi dans $\mathrm{cof} \cap W$ --- ce qui servira plus
bas. C'est le lemme de factorisation de K.~S.~Brown ; le nom n'est pas sur la
page.}

\emph{L'inverse proposé.} « Je dis que $(g, i)$ est l'inverse cherché. » Pour
voir que $(gf, i)$ est dans la composante de $\underline{1}_{x}$, on voudrait
utiliser le carré $\alpha i = f \cdot 1_{x}$ ; mais $\alpha$ n'est pas une
cofibration, et $\alpha i = f$, qui est dans $W$, ne l'est pas davantage.
Il faudrait une $\beta : y \to y'$ dans $W$ telle que $\beta f$ soit dans
$\mathrm{cof}$, ce qui n'est visiblement pas possible en général --- prendre
pour cofibrations les monomorphismes, et $f$ non mono. « Ça se présente
mal ! »

\emph{Le foncteur $q$.} On définit plutôt
\[
q : \widetilde{\mathcal{C}} \longrightarrow W^{-1}\mathcal{C},
\qquad x \mapsto x, \qquad (f, i) \mapsto [i]^{-1}[f],
\]
compatible aux identités et à la composition, et bien défini sur les
composantes : une flèche $u$ de $(f,i)$ vers $(f',i')$ est dans $W$ par
deux-sur-trois, et $[i']^{-1}[f'] = [ui]^{-1}[uf]$. Il est bijectif sur les
objets. Il est \emph{plein} : toute flèche de $W^{-1}\mathcal{C}$ est un
composé de $[f]$ et de $[w]^{-1}$, $w \in W$, et il suffit que $[f]^{-1}$
provienne de $\widetilde{\mathcal{C}}$ pour $f \in W$. Or $q(g, i) =
[i]^{-1}[g]$ et, de $\alpha i = f$, $[f][i]^{-1} = [\alpha]$, d'où
$[f]\,[i]^{-1}[g] = [\alpha][g] = [\alpha g] = [1_{y}]$.

\pagerange{11}{12}
\section*{Le doute, et la rectification (page 11)}

\emph{La difficulté} est qu'il n'est pas clair que $q$ soit \emph{fidèle} --- ce
qui en ferait un isomorphisme de catégories.

\emph{Le doute.} Le foncteur $q$ se factorise visiblement par
\[
\widetilde{\mathcal{C}} \longrightarrow W_{0}^{-1}\mathcal{C}
\longrightarrow W^{-1}\mathcal{C}, \qquad W_{0} = W \cap \mathrm{cof},
\]
et il faudrait que la seconde flèche soit une équivalence. « Mais je doute
qu'il ne suffise d'inverser les équiv.\ faibles \emph{cofibrantes}, pour que les
équiv.\ faibles deviennent inversibles ! »

Le doute n'est pas fondé, et la page 9 contient de quoi le lever. Quand tous
les objets sont cofibrants, toute $f \in W$ s'écrit $f = \alpha i$ avec
$\alpha g = 1_{y}$ et $i, g \in W_{0}$. Dans $W_{0}^{-1}\mathcal{C}$, $[g]$ est
inversible et $[\alpha][g] = 1$, donc $[\alpha] = [g]^{-1}$ est inversible, et
$[f] = [\alpha][i]$ aussi. Ainsi $W_{0}^{-1}\mathcal{C} \to W^{-1}\mathcal{C}$
est un isomorphisme de catégories.\footnote{L'argument est le nôtre, mais
chacun de ses éléments est sur la page 9. Il ne règle pas la question de la
page : que $W_{0}^{-1}\mathcal{C} = W^{-1}\mathcal{C}$ ne dit pas que les
fractions $[i]^{-1}[f]$ à un seul dénominateur, prises à composante connexe
près, suffisent à décrire les flèches, ni que deux fractions égales dans la
localisation soient dans la même composante. C'est la fidélité de $q$ qui
reste ouverte.}

\emph{La rectification.} Il change donc de définition : $\mathrm{Ch}'_{0}(x,y)$
est la catégorie des diagrammes
\[
x \xrightarrow{\;f\;} \widetilde{y} \xleftarrow{\;i\;} y,
\qquad f \in \mathrm{cof}, \quad i \in W,
\]
et $\widetilde{\mathcal{C}}'$ la catégorie des $\pi_{0}(\mathrm{Ch}'_{0}(x,y))$.
La page 5 demandait $i \in \mathrm{cof} \cap W$ sans rien demander à $f$ ;
la page 11 déplace la condition de cofibration du dénominateur au
numérateur.\footnote{Il écrit $\underline{Ch}_{0}$ et $\widetilde{\mathcal{C}}$
pour les deux constructions. C'est le point de bascule du dossier ; on les
distingue par un prime.}

\pagerange{13}{15}
\section*{Comparer deux factorisations, et le critère $(\star\star)$ (pages 13 à 15)}

\emph{L'image d'une flèche.} Soit $f : x \to y$. Par C3, $f = \alpha f_{0}$ avec
$f_{0} : x \to \bar{y}$ cofibration et $\alpha : \bar{y} \to y$ dans $W$ ; et,
$y$ étant cofibrant, on peut supposer que $\alpha$ admet une section $i$,
comme à la page 9. Le diagramme $(f_{0}, i)$ est dans $\mathrm{Ch}'_{0}(x,y)$,
et $q(f_{0}, i) = [i]^{-1}[f_{0}] = [\alpha][f_{0}] = [f]$.\footnote{La page 13
est la moins lisible du dossier : un palimpseste rapide, deux blocs annulés,
et une note marginale en biais qu'on ne restitue pas. On donne la
construction dans la forme que les pages 14 et 15 permettent de reconstituer.}

\emph{Deux factorisations} $f = \alpha f_{0} = \alpha' f'_{0}$ se comparent à
une troisième. La somme amalgamée de
\[
x \xleftarrow{\;\nabla_{x}\;} x \vee x \xrightarrow{\;f_{0} \vee f'_{0}\;}
\bar{y} \sqcup \bar{y}'
\]
existe, $f_{0} \vee f'_{0}$ étant une cofibration ; soit $\widetilde{y}$ cette
somme et $g : x \to \widetilde{y}$, cofibration par C2. Les flèches $(\alpha,
\alpha')$ et $f$ induisent $\varphi : \widetilde{y} \to y$, qu'on factorise en
$\varphi = \alpha''\varphi_{0}$, $\varphi_{0} \in \mathrm{cof}$,
$\alpha'' \in W$. Posant $f''_{0} = \varphi_{0}\, g$, on obtient des flèches de
factorisations $(f_{0}, \alpha) \to (f''_{0}, \alpha'')$ et
$(f'_{0}, \alpha') \to (f''_{0}, \alpha'')$.

\begin{tikzcd}[column sep=small, row sep=small, nodes={font=\scriptsize}]
x \vee x \arrow[r, "f_0 \vee f'_0"] \arrow[d, "\nabla_x"'] & \bar{y} \sqcup \bar{y}\,' \arrow[d] \arrow[dr, "{(\alpha, \alpha')}"] & \\
x \arrow[r, "g"'] & \widetilde{y} \arrow[r, "\varphi"'] & y
\end{tikzcd}

Il reste à comparer les sections\footnote{Le bloc de trois carrés de sommes amalgamées porté de sa
main dans la marge du listing de la page 14 est la préparation de ce
diagramme. Sur la page 15, $\varphi_{0}$ va de $\widetilde{y}$ dans un
$\bar{y}''$ et $\alpha''$ de $\bar{y}''$ dans $y$ ; le carré de gauche est
marqué « cocart ».} : il suffirait que les
sections de $\alpha$ et $\alpha'$ soient compatibles avec une section de
$\alpha''$, et ce n'est pas automatique.

\emph{La réduction.} Comme $(f_{0}, i) = (\mathrm{id}_{\bar{y}}, i) \circ
(f_{0}, \mathrm{id}_{\bar{y}})$, tout se ramène au choix de la section :

\begin{enumerate}
\item[$(\star\star)$] Soit $f : x \to y$ dans $W$, et $i, i' : y \to x$ deux
sections de $f$. Alors $(1_{x}, i)$ et $(1_{x}, i')$ sont dans la même
composante connexe de $\mathrm{Ch}'_{0}(x,y)$.
\end{enumerate}

\textbf{Énoncé final.} La condition $(\star\star)$ est nécessaire et suffisante
pour que $q' : \widetilde{\mathcal{C}}' \to W^{-1}\mathcal{C}$ soit fidèle,
donc un isomorphisme.

La nécessité est claire : $f i = 1_{y}$ donne $[i]^{-1} = [f]$, de sorte que
$q'(1_{x}, i) = [i]^{-1} = [f] = [i']^{-1} = q'(1_{x}, i')$ ; si $q'$ est
fidèle, les deux diagrammes définissent la même flèche de
$\widetilde{\mathcal{C}}'$, c'est-à-dire sont dans la même composante. La suffisance est affirmée, et le dossier s'arrête
là.\footnote{La page écrit « (nec.\ et) suffis. », la nécessité entre
parenthèses. Le nom de la construction dans la langue d'aujourd'hui est un
\emph{calcul de fractions à homotopie près}. P.~Gabriel et M.~Zisman (1967) ont
donné les conditions d'un calcul de fractions exact ; K.~S.~Brown (1973), pour
les catégories d'objets fibrants, et plus tard D.-C.~Cisinski, pour les
catégories dérivables, décrivent les flèches de la catégorie homotopique par
des fractions à un seul dénominateur prises à homotopie près, et non à
composante connexe près d'une catégorie de diagrammes. Cette lecture ne
tranche pas si $(\star\star)$ est vraie sous C1--C3 ; elle signale que la
littérature ne l'utilise pas sous cette forme, à notre connaissance.}

\end{document}
